\section{引言：为什么需要 RL 来做推理}

本讲由 CMU 的 Aviral Kumar 客座讲授，聚焦于强化学习（RL）在大语言模型（LLM）推理中的应用。讲者将推理限定在数学问题求解场景下进行阐述，并以时间线为轴将内容分为两大部分：经典 RL 技术和现代扩展。

\begin{importantbox}{核心动机}
传统的 next-token prediction（NTP）训练范式存在根本性局限：学习到的模型 $\hat{p}_\theta(\mathbf{y}|\mathbf{x})$ 与真实分布 $p^*(\mathbf{y}|\mathbf{x})$ 的误差正比于 $|\mathcal{D}(\mathbf{y}|\mathbf{x})|^{-\alpha}$，即误差随与目标输入 $\mathbf{x}$ 相似的数据量增加而减小。对于推理问题，高质量的问题--解答数据极其稀缺，预计到 2028 年互联网高质量文本数据将耗尽。
\end{importantbox}

\subsection{NTP 的失败模式}

仅靠监督训练的模型在解决如 AIME（美国数学邀请赛）或 IMO（国际数学奥林匹克）等困难数学问题时，虽然能生成"看上去像解答"的文本，但逻辑推理链条往往存在关键错误。模型会跳过必要的推理步骤，或者在证明过程中对个别项进行分析而忽略整体结构。

\begin{figure}[H]
\centering
\includegraphics[width=0.9\textwidth]{slides-images/slide-03.jpg}
\caption{传统训练方式：next-token prediction 的误差界\protect\footnotemark}
\end{figure}
\footnotetext{Slides 第 3 页。}

\begin{figure}[H]
\centering
\includegraphics[width=0.9\textwidth]{slides-images/slide-04.jpg}
\caption{数据稀缺问题：数学推理和具身 AI 领域的数据限制\protect\footnotemark}
\end{figure}
\footnotetext{Slides 第 4 页。}

\subsection{本章小结}

NTP 在推理场景下失效的核心原因是高质量推理数据不足。RL 的价值在于不依赖人工标注的完整解答，而是通过奖励信号（如答案是否正确）引导模型自主发现正确的推理路径。

